% Kết luận — lead viết ở Pha 3. Mọi con số lấy từ \SL{...} (outputs/metrics/*.json).
\chapter*{Kết luận}
\addcontentsline{toc}{chapter}{Kết luận}
\markboth{Kết luận}{Kết luận}
% chương không đánh số: hình/bảng đánh số riêng KL.1, KL.2... thay vì nối tiếp Chương 7
\setcounter{figure}{0}\setcounter{table}{0}
\renewcommand{\thefigure}{KL.\arabic{figure}}
\renewcommand{\thetable}{KL.\arabic{table}}

Tiểu luận đã đi một vòng từ lịch sử của trí tuệ nhân tạo đến các hệ thống thông minh nhiều thành phần. Mỗi chương
lý thuyết được kiểm chứng bằng thực nghiệm chạy thật trên máy cá nhân (GPU RTX~4060 8~GB và CPU 32 luồng), mã nguồn
in trong báo cáo được trích nguyên văn từ mã đã sinh ra số liệu, và mọi con số trong văn bản được đọc tự động từ
tệp kết quả. Phần này tổng hợp những gì đã làm được, rút ra các bài học xuyên suốt, nêu hạn chế và hướng phát triển.

\section*{Tổng hợp kết quả}
\addcontentsline{toc}{section}{Tổng hợp kết quả}

Chương~1 dựng lại tiến trình của AI qua các giai đoạn kỳ vọng và suy thoái, cho thấy mỗi lần phục hồi đều gắn với
một thay đổi về dữ liệu, phần cứng hoặc thuật toán. Các chương~2 đến~7 lần lượt hiện thực hóa từng họ phương pháp.
Bảng~\ref{tab:ch8-tongket} tóm tắt mô hình được chọn và chỉ số chính trên tập kiểm thử của từng thực nghiệm; mô hình
được chọn theo tập kiểm định (validation), không theo tập kiểm thử. Hình~\ref{fig:ch8-tongket} đặt các họ mô hình
cạnh dạng cấu trúc dữ liệu mà mỗi họ khai thác.

Toàn bộ mô hình của các chương~2 đến~7 được tích hợp vào một Web App Flask duy nhất. Mỗi chương cung cấp một mô-đun
phục vụ theo cùng một hợp đồng (danh sách thực nghiệm và hàm dự đoán), nên giao diện tự sinh biểu mẫu nhập liệu, nạp
mô hình khi được gọi lần đầu và hiển thị kết quả kèm chỉ số kiểm thử của mô hình. Trang chủ
(Hình~\ref{fig:ch8-home-web}) liệt kê 15 thực nghiệm theo chương; giao diện đáp ứng nên dùng được trên điện thoại
(Hình~\ref{fig:ch8-home-mobile}).

\begin{table}[!tbp]\centering
\caption{Tổng hợp kết quả các thực nghiệm (chỉ số trên tập kiểm thử)}\label{tab:ch8-tongket}
{\small\setlength{\tabcolsep}{4pt}
\begin{adjustbox}{max width=\linewidth}
\begin{tabular}{@{}llL{4.3cm}L{6.2cm}@{}}
\toprule
Ch. & Thực nghiệm & Mô hình được chọn & Chỉ số chính \\
\midrule
2 & Tiểu đường (100k) & \SL{ch2_diabetes/best/label} & F1 \SL{ch2_diabetes/best/f1}; ROC-AUC \SL{ch2_diabetes/best/auc} \\
2 & Giá nhà & \SL{ch2_house/best/label} & $R^2$ \SL{ch2_house/best/r2}; MAPE \SL{ch2_house/best/mapePct}\,\% \\
2 & Nhận xét khách hàng & \SL{ch2_comments/best/label} & macro-F1 \SL{ch2_comments/best/macro_f1}; Acc \SL{ch2_comments/best/accPct}\,\% \\
\midrule
3 & MNIST & Improved CNN (PyTorch) & Acc \SL{ch3_mnist/torch_improved/accPct}\,\% \\
3 & Fashion-MNIST & Improved CNN (PyTorch) & Acc \SL{ch3_fashion/torch_improved/accPct}\,\% \\
3 & Rock/Paper/Scissors & Improved CNN (PyTorch) & Acc \SL{ch3_rps/torch_improved/accPct}\,\% \\
3 & CIFAR-10 & Improved CNN (PyTorch) & Acc \SL{ch3_cifar10/torch_improved/accPct}\,\% \\
\midrule
4 & Doanh thu ngày & \SL{ch4_summary/ecommerce/best_arch} (\SL{ch4_summary/ecommerce/best_fw}) & RMSE bằng \SL{ch4_summary/ecommerce/best_ratioPct}\,\% của dự báo ngây thơ \\
4 & Cổ phiếu AAPL & \SL{ch4_summary/aapl/best_arch} (\SL{ch4_summary/aapl/best_fw}) & $R^2$ \SL{ch4_summary/aapl/best/r2}; RMSE bằng \SL{ch4_summary/aapl/best_ratioPct}\,\% dự báo ngây thơ \\
4 & Quỹ vàng GLD & \SL{ch4_summary/gold/best_arch} (\SL{ch4_summary/gold/best_fw}) & $R^2$ \SL{ch4_summary/gold/best/r2}; RMSE bằng \SL{ch4_summary/gold/best_ratioPct}\,\% dự báo ngây thơ \\
\midrule
5 & Truy xuất (ba miền) & Truy vấn KG; lai RRF + KG & MRR \SL{ch5_retrieval/sl/all/graph/mrrPct}\,\% và \SL{ch5_retrieval/sl/all/rrfg/mrrPct}\,\% (e5-small \SL{ch5_retrieval/sl/all/esmall/mrrPct}\,\%) \\
5 & Hỏi đáp cổ phiếu & Tác tử định tuyến công cụ & Acc \SL{ch5_agent_stock/systems/agent/accPct}\,\% (Advanced RAG \SL{ch5_agent_stock/systems/advanced_rag/accPct}\,\%, LLM thuần \SL{ch5_agent_stock/systems/llm_only/accPct}\,\%) \\
5 & Hỏi đáp giá vàng & Tác tử định tuyến công cụ & Acc \SL{ch5_agent_gold/systems/agent/accPct}\,\% (Advanced RAG \SL{ch5_agent_gold/systems/advanced_rag/accPct}\,\%, LLM thuần \SL{ch5_agent_gold/systems/llm_only/accPct}\,\%) \\
5 & Hỏi đáp tiểu đường & Tác tử định tuyến công cụ & Acc \SL{ch5_agent_diabetes/systems/agent/accPct}\,\% (Naive RAG \SL{ch5_agent_diabetes/systems/naive_rag/accPct}\,\%, LLM thuần \SL{ch5_agent_diabetes/systems/llm_only/accPct}\,\%) \\
\midrule
6 & Cora (phân loại nút) & \SL{ch6_summary/Cora/best} & Acc \SL{ch6_summary/Cora/best_accPct}\,\%; hơn MLP \SL{ch6_summary/Cora/gain_gat_mlpPct} điểm \\
6 & CiteSeer & \SL{ch6_summary/CiteSeer/best} & Acc \SL{ch6_summary/CiteSeer/best_accPct}\,\% \\
6 & Cora (dự đoán liên kết) & Graph Auto-Encoder & AUC \SL{ch6_link/test_auc_meanPct}\,\% \\
\midrule
7 & GSM8K (Qwen2.5-1.5B) & CoT một tác tử & Acc \SL{ch7_gsm8k/systems/cot/accPct}\,\% (đa tác tử \SL{ch7_gsm8k/systems/multi/accPct}\,\%) \\
\bottomrule
\end{tabular}
\end{adjustbox}}
\end{table}

\hinhtikz{ch8_tongket}{Bức tranh tổng thể: mỗi họ mô hình khai thác một dạng cấu trúc của dữ liệu, trên cùng một
chu trình dữ liệu -- baseline -- mô hình -- đánh giá -- triển khai}{ch8-tongket}

\anhUI{ch0_home_web}{Trang chủ Web App: 15 thực nghiệm của các chương~2--7 trong một ứng dụng}{ch8-home-web}

\anhMobile{ch0_home_mobile}{Trang chủ Web App trên khung điện thoại 390 px}{ch8-home-mobile}

\section*{Các bài học xuyên suốt}
\addcontentsline{toc}{section}{Các bài học xuyên suốt}

\begin{enumerate}
  \item \textbf{Cấu trúc của dữ liệu quyết định kiến trúc phù hợp.} Tích chập khai thác tính cục bộ của ảnh, hồi quy
  khai thác thứ tự thời gian, truyền thông điệp khai thác quan hệ trong đồ thị. Khi cấu trúc ấy bị phá, lợi thế biến
  mất: xáo trộn toàn bộ cạnh của Cora làm GCN chỉ còn \SL{ch6_summary/perturb/shuf100Pct}\,\%, thấp hơn cả MLP không
  dùng đồ thị; còn mạng GCN 8 tầng không có kết nối tắt rơi xuống \SL{ch6_summary/depth/acc8Pct}\,\% vì hiện tượng
  làm trơn quá mức (over-smoothing).
  \item \textbf{Baseline đơn giản là thước đo bắt buộc.} Trên dữ liệu giá nhà, Lasso tuyến tính thắng cả XGBoost và
  Random Forest vì quan hệ trong dữ liệu gần như tuyến tính. Với chuỗi giá cổ phiếu, mọi mô hình hồi quy chỉ ngang
  dự báo ``giữ nguyên giá hôm qua''; $R^2$ cao (\SL{ch4_summary/aapl/best/r2}) không có nghĩa là mô hình dự báo được
  biến động, vì độ chính xác hướng chỉ khoảng \SL{ch4_summary/aapl/best/dir_accPct}\,\%.
  \item \textbf{Cải tiến kiến trúc có hiệu quả tăng theo độ khó của bài toán.} Khối chuẩn hóa theo lô, dropout, tăng
  cường dữ liệu và kết nối tắt chỉ thêm \SL{ch3_mnist/gain_torch/accPct} điểm trên MNIST nhưng thêm
  \SL{ch3_cifar10/gain_torch/accPct} điểm trên CIFAR-10.
  \item \textbf{Tri thức có cấu trúc bù cho mô hình ngôn ngữ nhỏ.} Câu hỏi tra cứu chính xác (một ngày, một nhóm hồ
  sơ) được trả lời đúng bằng truy vấn đồ thị hoặc phép tính trên bảng, trong khi tìm kiếm thuần văn bản dễ nhầm các
  bản ghi giống nhau về con số. Trên bộ câu hỏi cổ phiếu, Qwen2.5-1.5B không có ngữ cảnh đúng
  \SL{ch5_agent_stock/systems/llm_only/accPct}\,\%, Naive RAG \SL{ch5_agent_stock/systems/naive_rag/accPct}\,\%,
  Advanced RAG \SL{ch5_agent_stock/systems/advanced_rag/accPct}\,\%, còn tác tử định tuyến sang công cụ tra KG và tính
  trên bảng đúng \SL{ch5_agent_stock/systems/agent/accPct}\,\%. Truy xuất tốt hơn cũng chưa đủ: trên miền tiểu đường,
  Advanced RAG (\SL{ch5_agent_diabetes/systems/advanced_rag/accPct}\,\%) còn thua Naive RAG
  (\SL{ch5_agent_diabetes/systems/naive_rag/accPct}\,\%) vì mô hình 1,5B chép sai con số dù tài liệu đúng đã nằm trong
  ngữ cảnh. Tác tử định tuyến công cụ là cách đơn giản và hiệu quả
  để kết hợp hai nguồn, với điều kiện câu hỏi thuộc các mẫu mà bộ định tuyến đã biết.
  \item \textbf{Nhiều tác tử hơn không tự động tốt hơn.} Hệ đa tác tử dùng trung bình \SL{ch7_gsm8k/systems/multi/calls}
  lượt gọi mô hình mỗi câu nhưng không vượt tác tử suy luận từng bước; giá trị thật của nó là tín hiệu tin cậy: khi
  hai bộ giải đồng ý, câu trả lời đúng \SL{ch7_gsm8k/multi/agree_accPct}\,\%.
  \item \textbf{Đánh giá trung thực quan trọng ngang mô hình.} Chia dữ liệu theo thời gian và chuẩn hóa chỉ trên tập
  huấn luyện tránh rò rỉ thông tin tương lai; chọn mô hình và ngưỡng trên tập kiểm định; báo cáo cả trường hợp mô hình
  không thắng baseline thay vì chỉ chọn kết quả đẹp.
\end{enumerate}

\section*{Hạn chế}
\addcontentsline{toc}{section}{Hạn chế}

\begin{itemize}
  \item Các thực nghiệm chạy trên một máy cá nhân, nhiều làn song song chia sẻ một GPU 8~GB, nên một số cấu hình
  phải rút gọn: mô hình TensorFlow chỉ chạy trên CPU (TensorFlow trên Windows không hỗ trợ GPU) nên chỉ được 6--7
  epoch và vẫn kém bản PyTorch cùng kiến trúc trên các bộ khó (Chương~3), cài đặt NumPy từ đầu
  dùng tập con, và thực nghiệm đa tác tử dùng tập con GSM8K (số câu của từng hệ được ghi rõ ở Chương~7).
  \item Phần lớn kết quả là một lần chạy với hạt giống cố định; chỉ Chương~6 báo cáo trung bình và độ lệch qua năm
  hạt giống. Khác biệt nhỏ giữa các mô hình vì vậy cần được đọc thận trọng.
  \item Mô hình ngôn ngữ dùng trong RAG và tác tử là Qwen2.5-1.5B chạy cục bộ; chất lượng diễn đạt tiếng Việt và khả
  năng suy luận của nó thấp hơn các mô hình lớn.
  \item Dữ liệu y tế và tài chính chỉ phục vụ mục đích học thuật; hệ thống không thay thế tư vấn y khoa hay khuyến
  nghị đầu tư.
\end{itemize}

\section*{Hướng phát triển}
\addcontentsline{toc}{section}{Hướng phát triển}

\begin{itemize}
  \item Lặp lại thực nghiệm với nhiều hạt giống và kiểm định thống kê để kết luận so sánh chắc chắn hơn.
  \item Dùng học chuyển giao (mô hình CNN tiền huấn luyện) cho các bộ ảnh nhỏ như Rock/Paper/Scissors, và mô hình
  Transformer cho chuỗi thời gian dài.
  \item Kết hợp GNN với đồ thị tri thức của Chương~5 (biểu diễn thực thể bằng KG embedding, GraphRAG có học).
  \item Thay mô hình ngôn ngữ 1,5B bằng mô hình lớn hơn hoặc tinh chỉnh nhẹ (LoRA) trên dữ liệu miền, và bổ sung bộ
  kiểm chứng trích dẫn tự động cho chatbot.
  \item Đóng gói Web App thành dịch vụ có giám sát (ghi log truy vấn, đo độ trễ, phát hiện trôi dữ liệu) để tiến gần
  một hệ thống thông minh vận hành thực tế.
\end{itemize}
